\chapter{Reprodução dos experimentos}\label{ap:experimentos}

Todas as tabelas e figuras do \autoref{cap:resultados} que não são capturas de tela foram geradas pelos roteiros em Python da pasta \arquivo{artigo/experimentos/}. Eles reimplementam, em menos de duzentas linhas, o núcleo físico do simulador: a força da \autoref{eq:nbody}, os integradores leapfrog, Yoshida e RK4, os elementos osculadores, a equação de Kepler e a série de Fourier. Leem as mesmas condições iniciais do programa.

A \autoref{tab:validacao} gerada por esses roteiros coincide, até a última casa publicada, com a tabela de validação que a documentação da versão 2D atribui ao próprio programa. Como o algoritmo e o passo são os mesmos, isso confirma o porte para outra linguagem; a validação da física é a comparação com o JPL.

\begin{table}[htb]
  \centering
  \caption{Roteiros de reprodução}
  \label{tab:roteiros}
  \small
  \begin{tabularx}{\textwidth}{lX}
    \toprule
    Arquivo & O que produz \\
    \midrule
    \arquivo{nbody.py} & força, integradores, energia e momento angular \\
    \arquivo{fourier\_lib.py} & elementos osculadores, equação de Kepler, DFT e forma fechada de Bessel \\
    \arquivo{e1\_validacao.py} & \autoref{tab:validacao} e \autoref{fig:validacao} \\
    \arquivo{e2\_integradores.py} & dados das \autoref{fig:trabalho-precisao} e \ref{fig:energia} \\
    \arquivo{e4\_copernico.py} & \autoref{tab:copernico-fourier} \\
    \arquivo{figuras.py} & figuras dos \autoref{cap:matematica} e \ref{cap:resultados} \\
    \arquivo{figuras2.py} & \autoref{fig:angulo-tempo} \\
    \arquivo{figuras3.py} & \autoref{fig:marte-retrogrado} e \autoref{fig:sol-baricentro} \\
    \arquivo{fig\_elementos.py} & \autoref{fig:elementos} \\
    \bottomrule
  \end{tabularx}
  \legend{Requer Python 3 com NumPy, SciPy e Matplotlib. Instruções em \arquivo{experimentos/LEIAME.md}. Fonte: elaborada pelos autores.}
\end{table}

\lstinputlisting[language=Python,firstline=1,lastline=60,
  caption={Núcleo da reimplementação de referência (\arquivo{experimentos/nbody.py})},label=cod:nbody-py]{experimentos/nbody.py}
